% GENERATED FILE. Edit canonical lessons, then run npm run generate:books.
% canonical-source-hash: fnv1a64:0028235d59dad308
% canonical-lessons: BN-C66-pahar, BN-C66-bon, BN-C66-nodi, BN-C66-pukur, BN-C66-pathor

\chapter{Mountain, Forest, River, Pond, Stone}
\label{ch:bn-mountain-forest-river-pond-stone}
\hlchaptermodality{66}

\begin{chapteropening}
\textbf{By the end of this chapter:} \emph{I can say a mountain, a forest, a river, a pond and a stone.}

Complete the last lesson of chapter 66: i can say a mountain, a forest, a river, a pond and a stone.
\end{chapteropening}

\section[pāhāṛ]{\bn{পাহাড়} (pāhāṛ) — a mountain}

\label{lesson:BN-C66-pahar}



\begin{quote}
Before the new one: say the Bengali for a flower, then the Bengali for a leaf.
\end{quote}

\subsection*{You'll want to know: \bn{পাহাড়}}
\textbf{\bn{পাহাড়}} — \emph{pāhāṛ} — \textquotedblleft{}a mountain\textquotedblright{}.

\begin{grammarlens}[title={What you've built}]
One more word for this chapter.
\end{grammarlens}

\subsection*{Your turn}
\begin{itemize}
\raggedright
  \item \emph{Say it:} \emph{pāhāṛ}
  \item \emph{Say it:} \emph{pāhāṛ}, once more
  \item \emph{Say it:} say \emph{pāhāṛ}
  \item \emph{Recall:} say the Bengali for a flower, then the Bengali for a leaf, then say \emph{pāhāṛ} again
\end{itemize}

\begin{tcolorbox}[breakable,colback=teal!4,colframe=teal!35!black,title={Before you move on}]
What does \bn{পাহাড়} mean? (\textquotedblleft{}A mountain\textquotedblright{}.) Say it once more.
\end{tcolorbox}
\section[bôn]{\bn{বন} (bôn) — a forest}

\label{lesson:BN-C66-bon}



\begin{quote}
Before the new one: say the Bengali for a leaf, then the Bengali for a mountain.
\end{quote}

\subsection*{You'll want to know: \bn{বন}}
\textbf{\bn{বন}} — \emph{bôn} — \textquotedblleft{}a forest\textquotedblright{}.

\begin{grammarlens}[title={What you've built}]
One more word for this chapter.
\end{grammarlens}

\subsection*{Your turn}
\begin{itemize}
\raggedright
  \item \emph{Say it:} \emph{bôn}
  \item \emph{Say it:} \emph{bôn}, once more
  \item \emph{Say it:} say \emph{bôn}
  \item \emph{Recall:} say the Bengali for a leaf, then the Bengali for a mountain, then say \emph{bôn} again
\end{itemize}

\begin{tcolorbox}[breakable,colback=teal!4,colframe=teal!35!black,title={Before you move on}]
What does \bn{বন} mean? (\textquotedblleft{}A forest\textquotedblright{}.) Say it once more.
\end{tcolorbox}
\section[nôdi]{\bn{নদী} (nôdi) — a river}

\label{lesson:BN-C66-nodi}



\begin{quote}
Before the new one: say the Bengali for a mountain, then the Bengali for a forest.
\end{quote}

\subsection*{You'll want to know: \bn{নদী}}
\textbf{\bn{নদী}} — \emph{nôdi} — \textquotedblleft{}a river\textquotedblright{}.

\begin{grammarlens}[title={What you've built}]
One more word for this chapter.
\end{grammarlens}

\subsection*{Your turn}
\begin{itemize}
\raggedright
  \item \emph{Say it:} \emph{nôdi}
  \item \emph{Say it:} \emph{nôdi}, once more
  \item \emph{Say it:} say \emph{nôdi}
  \item \emph{Recall:} say the Bengali for a mountain, then the Bengali for a forest, then say \emph{nôdi} again
\end{itemize}

\begin{tcolorbox}[breakable,colback=teal!4,colframe=teal!35!black,title={Before you move on}]
What does \bn{নদী} mean? (\textquotedblleft{}A river\textquotedblright{}.) Say it once more.
\end{tcolorbox}
\section[pukur]{\bn{পুকুর} (pukur) — a pond}

\label{lesson:BN-C66-pukur}



\begin{quote}
Before the new one: say the Bengali for a forest, then the Bengali for a river.
\end{quote}

\subsection*{You'll want to know: \bn{পুকুর}}
\textbf{\bn{পুকুর}} — \emph{pukur} — \textquotedblleft{}a pond\textquotedblright{}.

\begin{grammarlens}[title={What you've built}]
One more word for this chapter.
\end{grammarlens}

\subsection*{Your turn}
\begin{itemize}
\raggedright
  \item \emph{Say it:} \emph{pukur}
  \item \emph{Say it:} \emph{pukur}, once more
  \item \emph{Say it:} say \emph{pukur}
  \item \emph{Recall:} say the Bengali for a forest, then the Bengali for a river, then say \emph{pukur} again
\end{itemize}

\begin{tcolorbox}[breakable,colback=teal!4,colframe=teal!35!black,title={Before you move on}]
What does \bn{পুকুর} mean? (\textquotedblleft{}A pond\textquotedblright{}.) Say it once more.
\end{tcolorbox}
\section[pāthôr]{\bn{পাথর} (pāthôr) — a stone}

\label{lesson:BN-C66-pathor}



\begin{quote}
Before the new one: say the Bengali for a river, then the Bengali for a pond.
\end{quote}

\subsection*{You'll want to know: \bn{পাথর}}
\textbf{\bn{পাথর}} — \emph{pāthôr} — \textquotedblleft{}a stone\textquotedblright{}.

\begin{grammarlens}[title={What you've built}]
One more word for this chapter.
\end{grammarlens}

\subsection*{Your turn}
\begin{itemize}
\raggedright
  \item \emph{Say it:} \emph{pāthôr}
  \item \emph{Say it:} \emph{pāthôr}, once more
  \item \emph{Say it:} say \emph{pāthôr}
  \item \emph{Recall:} say the Bengali for a river, then the Bengali for a pond, then say \emph{pāthôr} again
\end{itemize}

\begin{tcolorbox}[breakable,colback=teal!4,colframe=teal!35!black,title={Before you move on}]
What does \bn{পাথর} mean? (\textquotedblleft{}A stone\textquotedblright{}.) Say it once more.
\end{tcolorbox}
